\chapter{Multilinear Algebra}\label{multilinear-algebra}
\chapterstatus{complete}

\section{Introduction}\label{multilinear-algebra:section-introduction}

This chapter constructs the tensor, symmetric and exterior algebras of a vector space, introduces the
tensors of type $(r,s)$ of geometry and physics with their components, contractions and the raising and
lowering of indices, and studies the structures that underlie differential forms: alternating forms and interior products,
the Pfaffian of a skew-symmetric matrix, orientations, the Hodge star operator of an inner product space, and the decomposition of the
exterior algebra of a real vector space with a complex structure into forms of type $(p,q)$. These are
the pointwise versions of the constructions of Chapters~\ref{differential-forms},
\ref{riemannian-geometry} and~\ref{complex-manifolds}.

References are \cite[Chapters XVI and XIX]{lang-algebra} and \cite{greub-multilinear}.

\noindent\textbf{Prerequisites.} Chapter~\ref{rings} (Rings and Modules) and
Chapter~\ref{linear-algebra} (Linear Algebra).

\noindent\textbf{On the way to the Hodge conjecture.} The whole chapter is needed: tensor products, tensors of type $(r,s)$ and contraction, the exterior algebra, alternating forms, the Hodge star and forms of type $(p,q)$ are the
algebra underlying differential forms and the Hodge decomposition.

Throughout, $k$ is a field, vector spaces are over $k$ and $\otimes = \otimes_k$. Tensor products were
defined in Section~\ref{rings:section-tensor}.

\section{Tensor products of vector spaces}\label{multilinear-algebra:section-tensor-products}

\begin{proposition}\label{multilinear-algebra:proposition-tensor-basis}
Let $(v_i)_{i \in I}$ be a basis of $V$ and $(w_j)_{j \in J}$ a basis of $W$. Then $(v_i \otimes w_j)$ is a basis of
$V \otimes W$; in particular $\dim(V \otimes W) = \dim V \cdot \dim W$. For vector spaces
$V_1, \ldots, V_n, U$, multilinear maps $V_1 \times \cdots \times V_n \to U$ correspond bijectively to linear maps
$V_1 \otimes \cdots \otimes V_n \to U$.
\end{proposition}

\begin{proof}
$V \cong k^{(I)}$ and $W \cong k^{(J)}$, and by Proposition~\ref{rings:proposition-tensor-properties}
$k^{(I)} \otimes k^{(J)} \cong \bigoplus_{i} (k \otimes k^{(J)}) \cong k^{(I \times J)}$, with $e_i \otimes e_j \mapsto e_{(i,j)}$. The second
statement follows by induction from the definition of the tensor product and the associativity
isomorphism.
\end{proof}

\begin{proposition}\label{multilinear-algebra:proposition-hom-tensor}
The linear map $V^* \otimes W \to \Hom_k(V, W)$, $\varphi \otimes w \mapsto (v \mapsto \varphi(v)w)$, is injective, with image the
maps of finite rank; it is an isomorphism if $V$ or $W$ is finite-dimensional. For finite-dimensional
$V$, $W$ the map $V^* \otimes W^* \to (V \otimes W)^*$, $\varphi \otimes \psi \mapsto (v \otimes w \mapsto \varphi(v)\psi(w))$, is an
isomorphism. Under $V^* \otimes V \cong \End(V)$, the \emph{contraction} $\varphi \otimes v \mapsto \varphi(v)$ is the trace.
\end{proposition}

\begin{proof}
Let $(w_j)$ be a basis of $W$. An element of $V^* \otimes W$ is uniquely $\sum_j \varphi_j \otimes w_j$ with finitely
many nonzero $\varphi_j$ (Proposition~\ref{multilinear-algebra:proposition-tensor-basis} applied to a
basis of $V^*$); it maps to $v \mapsto \sum_j \varphi_j(v)w_j$, which is zero only if all $\varphi_j = 0$. The image
consists of maps with image in a finite-dimensional subspace, and every such map is of this form
(choose a basis of $W$ extending a basis of a finite-dimensional subspace containing the image). If $V$
is finite-dimensional, every map has finite rank; if $W$ is, likewise. For the second isomorphism,
compare dual bases. Finally, with a basis $(e_i)$ of $V$, $e_i^* \otimes e_j$ corresponds to the matrix unit
$E_{ji}$, whose trace is $\delta_{ij} = e_i^*(e_j)$.
\end{proof}

\begin{counterexample}\label{multilinear-algebra:counterexample-indecomposable}
Not every element of $V \otimes W$ is of the form $v \otimes w$. Under $k^2 \otimes k^2 \cong M_2(k)$ (sending
$e_i \otimes e_j$ to the matrix unit $E_{ij}$), the element $v \otimes w$ corresponds to a matrix of rank at most
$1$, while $e_1 \otimes e_1 + e_2 \otimes e_2$ corresponds to the identity matrix, of rank $2$.
\end{counterexample}

\section{The tensor, symmetric and exterior algebras}\label{multilinear-algebra:section-algebras}

\begin{definition}\label{multilinear-algebra:definition-tensor-algebra}
The \emph{tensor algebra} of $V$ is the graded $k$-algebra $T(V) = \bigoplus_{n \geq 0} T^n(V)$, with
$T^0(V) = k$, $T^n(V) = V^{\otimes n}$, and multiplication induced by
$(v_1 \otimes \cdots \otimes v_p)(w_1 \otimes \cdots \otimes w_q) = v_1 \otimes \cdots \otimes v_p \otimes w_1 \otimes \cdots \otimes w_q$. The \emph{symmetric
algebra} is $S(V) = T(V)/I_S$, where $I_S$ is the two-sided ideal generated by the elements
$v \otimes w - w \otimes v$, and the \emph{exterior algebra} is $\Lambda(V) = T(V)/I_\Lambda$, where $I_\Lambda$ is
generated by the elements $v \otimes v$. The ideals are homogeneous, so both quotients are graded:
$S(V) = \bigoplus_n S^n(V)$ and $\Lambda(V) = \bigoplus_n \Lambda^n(V)$. We write $v_1 \cdots v_n$ for the image of
$v_1 \otimes \cdots \otimes v_n$ in $S(V)$ and $v_1 \wedge \cdots \wedge v_n$ for its image in $\Lambda(V)$.
\end{definition}

\begin{proposition}\label{multilinear-algebra:proposition-universal-algebras}
Let $A$ be an associative $k$-algebra and $f \colon V \to A$ linear.
\begin{enumerate}
\item $f$ extends uniquely to an algebra homomorphism $T(V) \to A$.
\item If $A$ is commutative, it extends uniquely to $S(V) \to A$.
\item If $f(v)^2 = 0$ for all $v \in V$, it extends uniquely to $\Lambda(V) \to A$.
\end{enumerate}
Moreover, $n$-linear maps $V^n \to U$ that are symmetric (respectively alternating) correspond to linear maps
$S^n(V) \to U$ (respectively $\Lambda^n(V) \to U$).
\end{proposition}

\begin{proof}
(1) The $n$-linear map $(v_1, \ldots, v_n) \mapsto f(v_1) \cdots f(v_n)$ induces $T^n(V) \to A$, and the sum of these is an
algebra homomorphism; it is unique since $V$ generates $T(V)$. (2), (3) The homomorphism of (1) kills the
generators of $I_S$, respectively $I_\Lambda$.

For the last statement, the degree-$n$ part of $I_\Lambda$ is spanned by the tensors
$x \otimes v \otimes v \otimes y$ with $x, y$ pure tensors, and similarly for $I_S$. A multilinear $\beta \colon V^n \to U$
vanishes on the tensors with two equal adjacent factors if and only if it is alternating: expanding
$\beta(\ldots, v + w, v + w, \ldots) = 0$ shows it changes sign under adjacent transpositions, hence under all
transpositions up to sign, so it vanishes whenever any two arguments are equal. The symmetric case is
analogous.
\end{proof}

\begin{proposition}\label{multilinear-algebra:proposition-exterior-properties}
\begin{enumerate}
\item In $\Lambda(V)$, $v \wedge w = -w \wedge v$ for $v, w \in V$, and $\alpha \wedge \beta = (-1)^{pq}\beta \wedge \alpha$ for
$\alpha \in \Lambda^p(V)$, $\beta \in \Lambda^q(V)$.
\item $S(V)$ is commutative. If $(e_1, \ldots, e_n)$ is a basis of $V$, then $S(V) \cong k[x_1, \ldots, x_n]$ with
$e_i \mapsto x_i$, and $S^d(V)$ corresponds to the homogeneous polynomials of degree $d$.
\item If $(e_1, \ldots, e_n)$ is a basis of $V$, the elements $e_I = e_{i_1} \wedge \cdots \wedge e_{i_p}$, for
$I = \{i_1 < \cdots < i_p\} \subseteq \{1, \ldots, n\}$, form a basis of $\Lambda^p(V)$. So
$\dim \Lambda^p(V) = \binom{n}{p}$, $\Lambda^p(V) = 0$ for $p > n$, and $\dim\Lambda(V) = 2^n$.
\item Vectors $v_1, \ldots, v_p$ are linearly independent if and only if $v_1 \wedge \cdots \wedge v_p \neq 0$.
\item For $f \colon V \to W$ linear there are induced maps $\Lambda^p f \colon \Lambda^p V \to \Lambda^p W$,
$v_1 \wedge \cdots \wedge v_p \mapsto f(v_1) \wedge \cdots \wedge f(v_p)$, and if $\dim V = n$, $\Lambda^n f$ is multiplication by
$\det f$ for $f \colon V \to V$.
\end{enumerate}
\end{proposition}

\begin{proof}
(1) $0 = (v + w) \wedge (v + w) = v \wedge w + w \wedge v$. Moving each of the $q$ factors of a decomposable $\beta$
past the $p$ factors of $\alpha$ gives $pq$ sign changes.

(2) Commutativity holds on generators. The two universal properties
(Proposition~\ref{multilinear-algebra:proposition-universal-algebras}(2) and
Proposition~\ref{rings:proposition-polynomial-universal}) give inverse homomorphisms.

(3) By (1) and $e_i \wedge e_i = 0$, the $e_I$ span $\Lambda^p(V)$. For independence, fix $I$ and consider the
map $V^p \to k$, $(v_1, \ldots, v_p) \mapsto \det(e_{i_a}^*(v_b))_{a,b}$. It is multilinear and alternating by
Theorem~\ref{linear-algebra:theorem-determinant}, so it induces $\varphi_I \colon \Lambda^p(V) \to k$, and
$\varphi_I(e_J) = \delta_{IJ}$. A relation $\sum_J c_J e_J = 0$ gives $c_I = 0$ after applying $\varphi_I$.

(4) If $v_1 = \sum_{i \geq 2} c_i v_i$, the product vanishes by (1) and $v_i \wedge v_i = 0$. If they are independent,
extend them to a basis and use (3).

(5) The map exists by the universal property. For $V = W$ of dimension $n$ with basis $(e_i)$,
$(v_1, \ldots, v_n) \mapsto$ the coefficient of $f(v_1) \wedge \cdots \wedge f(v_n)$ on $e_1 \wedge \cdots \wedge e_n$ is multilinear
alternating; by Theorem~\ref{linear-algebra:theorem-determinant} applied to the columns of the matrix
of $f$, $f(e_1) \wedge \cdots \wedge f(e_n) = \det(f)\, e_1 \wedge \cdots \wedge e_n$.
\end{proof}

\begin{proposition}\label{multilinear-algebra:proposition-minors}
Let $f \colon V \to W$ be linear with matrix $A = (a_{ij})$ in bases $(e_1, \ldots, e_m)$ of $V$ and $(e'_1, \ldots, e'_n)$ of $W$. For subsets $I \subseteq \{1, \ldots, n\}$, $J \subseteq \{1, \ldots, m\}$ of size $p$, let
$A_{I,J}$ be the $p \times p$ submatrix with rows $I$ and columns $J$.
\begin{enumerate}
\item The matrix of $\Lambda^pf$ in the bases $(e_J)$ and $(e'_I)$ is $(\det A_{I,J})_{I,J}$: $\Lambda^pf(e_J) = \sum_I\det(A_{I,J})\,e'_I$.
\item (Cauchy--Binet) For an $n \times m$ matrix $A$ and an $m \times n$ matrix $B$, $\det(AB) = \sum_S\det A_{\{1..n\},S}\det B_{S,\{1..n\}}$, the sum over the subsets $S \subseteq \{1, \ldots, m\}$ of size $n$; in
particular $\det(AB) = 0$ if $m < n$, and $\det(AB) = \det A\det B$ if $m = n$.
\item The rank of $f$ is the largest $p$ with $\Lambda^pf \neq 0$, that is, the largest size of a nonzero minor of $A$.
\end{enumerate}
\end{proposition}

\begin{proof}
(1) $\Lambda^pf(e_J) = f(e_{j_1}) \wedge \cdots \wedge f(e_{j_p})$, and its coefficient on $e'_I$ is obtained by applying the functional $\varphi_I$ of the proof of
Proposition~\ref{multilinear-algebra:proposition-exterior-properties}(3), which gives $\det(e'^*_{i_a}(f(e_{j_b})))_{a,b} = \det(a_{i_aj_b}) = \det A_{I,J}$.

(2) By functoriality $\Lambda^n(AB) = \Lambda^nA \circ \Lambda^nB$ on $\Lambda^n(k^n)$, which is one-dimensional with basis $e_{\{1..n\}}$, and $\Lambda^n(AB)$ is multiplication by $\det(AB)$
(Proposition~\ref{multilinear-algebra:proposition-exterior-properties}(5)). By (1), $\Lambda^nB(e_{\{1..n\}}) = \sum_S\det B_{S,\{1..n\}}e_S$ and $\Lambda^nA(e_S) = \det A_{\{1..n\},S}\,e_{\{1..n\}}$. If $m < n$ there are
no such $S$.

(3) The image of $\Lambda^pf$ is spanned by the products $f(v_1) \wedge \cdots \wedge f(v_p)$, so it is the image of $\Lambda^p(\operatorname{im}f) \to \Lambda^pW$, which is injective because a basis of $\operatorname{im}f$ extends
to a basis of $W$ and Proposition~\ref{multilinear-algebra:proposition-exterior-properties}(3) applies to both. So $\Lambda^pf \neq 0$ if and only if $\Lambda^p(\operatorname{im}f) \neq 0$, that is, $p \leq \operatorname{rank}f$.
Combine with (1).
\end{proof}

\begin{counterexample}\label{multilinear-algebra:counterexample-characteristic-two}
In characteristic $2$ symmetric and alternating notions separate. Let $k = \mathbb{F}_2$ and $V = k$. The bilinear form $B(x, y) = xy$ satisfies $B(y, x) = B(x, y) = -B(x, y)$, so it is both symmetric and
antisymmetric, but it is not alternating, since $B(1, 1) = 1$; so it does not factor through $\Lambda^2V = 0$. Similarly the symmetrisation $S^2V \to V \otimes V$, $vw \mapsto v \otimes w + w \otimes v$, sends the generator $e^2$
of $S^2V$ to $2\,e \otimes e = 0$, so it is not injective. Factorials cause the same phenomena elsewhere: in Proposition~\ref{multilinear-algebra:proposition-symplectic-determinant} the characteristic enters through the factor $m!$ in
$\omega^m = m!\,e_1^* \wedge f_1^* \wedge \cdots \wedge e_m^* \wedge f_m^*$, which vanishes in characteristic $p \leq m$.
\end{counterexample}

\section{Tensors of type $(r,s)$}\label{multilinear-algebra:section-tensors}

The word \emph{tensor} is used in two related ways. In algebra a tensor is any element of a tensor product. In geometry and physics a tensor is a multilinear
quantity attached to a single vector space $V$ (in geometry, the tangent space at a point), with a certain number of slots for vectors and a certain number of
slots for linear forms: an inner product eats two vectors and returns a number, an endomorphism eats a vector and returns a vector, the curvature of a Riemannian
manifold eats three vectors and returns a vector. Older texts, and most physics texts, take yet another point of view: a tensor is an array of numbers, one for each
choice of basis, which changes in a prescribed way when the basis is changed. This section explains why the three points of view agree, and introduces the operations
on tensors (tensor product, contraction, raising and lowering of indices) that are used constantly in differential geometry.

In this section $V$ has finite dimension $n$, and $(e_1, \ldots, e_n)$ is a basis of $V$. Following the custom of differential geometry, we write
$e^i = e_i^*$ for the dual basis of $V^*$ (Definition~\ref{linear-algebra:definition-dual}), so that $e^i(e_j) = \delta^i_j$, where $\delta^i_j$ is $1$ if $i = j$
and $0$ otherwise. Indices of basis vectors of $V$ are written below and indices of basis vectors of $V^*$ above; the reason for this convention will appear
in Proposition~\ref{multilinear-algebra:proposition-tensor-transformation}.

\begin{definition}\label{multilinear-algebra:definition-tensor-type}
Let $r, s \in \NN$. A \emph{tensor of type $(r,s)$} on $V$ is an element of
\[
T^r_s(V) = \underbrace{V \otimes \cdots \otimes V}_{r} \otimes \underbrace{V^* \otimes \cdots \otimes V^*}_{s} = V^{\otimes r} \otimes (V^*)^{\otimes s},
\]
with the convention $T^0_0(V) = k$. One says that such a tensor is \emph{$r$ times contravariant} and \emph{$s$ times covariant}. The \emph{tensor product} of
$S = v_1 \otimes \cdots \otimes v_r \otimes \varphi^1 \otimes \cdots \otimes \varphi^s \in T^r_s(V)$ and $S^{\prime} = v^{\prime}_1 \otimes \cdots \otimes v^{\prime}_{r^{\prime}}
\otimes \psi^1 \otimes \cdots \otimes \psi^{s^{\prime}} \in T^{r^{\prime}}_{s^{\prime}}(V)$ is
\[
S \otimes S^{\prime} = v_1 \otimes \cdots \otimes v_r \otimes v^{\prime}_1 \otimes \cdots \otimes v^{\prime}_{r^{\prime}} \otimes \varphi^1 \otimes \cdots \otimes \varphi^s \otimes \psi^1 \otimes \cdots
\otimes \psi^{s^{\prime}} \in T^{r + r^{\prime}}_{s + s^{\prime}}(V),
\]
extended bilinearly; that is, it is the tensor product of Section~\ref{multilinear-algebra:section-tensor-products} followed by the isomorphism that moves the
factors $V$ in front of the factors $V^*$.
\end{definition}

The tensor product of tensors is well defined: the expression on the right is multilinear in the $r + s + r^{\prime} + s^{\prime}$ vectors and forms, so it induces a
bilinear map $T^r_s(V) \times T^{r^{\prime}}_{s^{\prime}}(V) \to T^{r + r^{\prime}}_{s + s^{\prime}}(V)$ by Proposition~\ref{multilinear-algebra:proposition-tensor-basis}. It is associative,
and it is not commutative: for $v, w \in V$ independent, $v \otimes w \neq w \otimes v$ in $T^2_0(V)$ (Proposition~\ref{multilinear-algebra:proposition-tensor-basis}).
The small cases are familiar: $T^0_0(V) = k$ consists of scalars, $T^1_0(V) = V$ of vectors and $T^0_1(V) = V^*$ of linear forms, also called \emph{covectors}.
The next proposition identifies the others with multilinear maps.

\begin{proposition}\label{multilinear-algebra:proposition-tensors-multilinear}
\begin{enumerate}
\item There is a unique isomorphism from $T^r_s(V)$ to the space of multilinear maps $(V^*)^r \times V^s \to k$ sending
$v_1 \otimes \cdots \otimes v_r \otimes \varphi^1 \otimes \cdots \otimes \varphi^s$ to the map
\[
(\psi^1, \ldots, \psi^r, w_1, \ldots, w_s) \longmapsto \psi^1(v_1) \cdots \psi^r(v_r)\,\varphi^1(w_1) \cdots \varphi^s(w_s).
\]
\item There is a unique isomorphism from $T^1_s(V)$ to the space of multilinear maps $V^s \to V$ sending $v \otimes \varphi^1 \otimes \cdots \otimes \varphi^s$ to the map
$(w_1, \ldots, w_s) \mapsto \varphi^1(w_1) \cdots \varphi^s(w_s)\,v$. For $s = 1$ this is the isomorphism $T^1_1(V) \cong \End(V)$ of
Proposition~\ref{multilinear-algebra:proposition-hom-tensor}.
\item The tensors $e_{i_1} \otimes \cdots \otimes e_{i_r} \otimes e^{j_1} \otimes \cdots \otimes e^{j_s}$, for all $i_1, \ldots, i_r, j_1, \ldots, j_s \in \{1, \ldots, n\}$, form a basis of
$T^r_s(V)$. In particular $\dim T^r_s(V) = n^{r+s}$.
\end{enumerate}
\end{proposition}

\begin{proof}
(3) The $e^j$ form a basis of $V^*$ (Proposition~\ref{linear-algebra:proposition-duality}), and Proposition~\ref{multilinear-algebra:proposition-tensor-basis}, applied
$r + s - 1$ times, says that the tensor products of basis vectors of the factors form a basis of the tensor product.

(1) By Proposition~\ref{multilinear-algebra:proposition-tensor-basis}, multilinear maps $(V^*)^r \times V^s \to k$ are the same as linear maps
$(V^*)^{\otimes r} \otimes V^{\otimes s} \to k$, that is, elements of $U^*$ with $U = (V^*)^{\otimes r} \otimes V^{\otimes s}$. Both $T^r_s(V)$ and $U^*$ have dimension $n^{r+s}$,
by (3) and Proposition~\ref{linear-algebra:proposition-duality}. The prescribed map exists, because the displayed expression is multilinear in
$(v_1, \ldots, v_r, \varphi^1, \ldots, \varphi^s)$, and it is unique, because the decomposable tensors span $T^r_s(V)$. It sends the basis vector
$e_{i_1} \otimes \cdots \otimes e_{i_r} \otimes e^{j_1} \otimes \cdots \otimes e^{j_s}$ to the linear form on $U$ that takes the value $1$ on
$e^{i_1} \otimes \cdots \otimes e^{i_r} \otimes e_{j_1} \otimes \cdots \otimes e_{j_s}$ and $0$ on the other basis vectors of $U$ of the same kind. These linear forms are the dual basis of the
basis of $U$ given by (3) (with the roles of $V$ and $V^*$ exchanged), so the map sends a basis to a basis and is an isomorphism.

(2) The same argument applies: multilinear maps $V^s \to V$ are linear maps $V^{\otimes s} \to V$, a space of dimension $n^s \cdot n = \dim T^1_s(V)$; the prescribed map
exists and is unique as in (1), and it sends $e_i \otimes e^{j_1} \otimes \cdots \otimes e^{j_s}$ to the linear map $V^{\otimes s} \to V$ that sends
$e_{j_1} \otimes \cdots \otimes e_{j_s}$ to $e_i$ and the other basis vectors $e_{l_1} \otimes \cdots \otimes e_{l_s}$ to $0$. These maps form a basis of $\Hom(V^{\otimes s}, V)$ (they are the
matrix units in the bases given by (3)), so the map is an isomorphism. For $s = 1$ it sends $v \otimes \varphi$ to $w \mapsto \varphi(w)v$, which is the map of
Proposition~\ref{multilinear-algebra:proposition-hom-tensor} composed with the exchange of factors $V \otimes V^* \cong V^* \otimes V$.
\end{proof}

From now on we identify a tensor of type $(r,s)$ with the corresponding multilinear map of (1), and a tensor of type $(1,s)$ also with the multilinear map
$V^s \to V$ of (2); the two are related by $T(\psi, w_1, \ldots, w_s) = \psi(T(w_1, \ldots, w_s))$, as one checks on decomposable tensors. So a bilinear form on $V$ is
the same as a tensor of type $(0,2)$, and an endomorphism the same as a tensor of type $(1,1)$.

\begin{definition}\label{multilinear-algebra:definition-tensor-components}
The \emph{components} of $T \in T^r_s(V)$ in the basis $(e_i)$ are the numbers
\[
T^{i_1 \ldots i_r}_{j_1 \ldots j_s} = T(e^{i_1}, \ldots, e^{i_r}, e_{j_1}, \ldots, e_{j_s}),
\]
so that, by Proposition~\ref{multilinear-algebra:proposition-tensors-multilinear},
\[
T = \sum_{i_1, \ldots, i_r, j_1, \ldots, j_s} T^{i_1 \ldots i_r}_{j_1 \ldots j_s}\; e_{i_1} \otimes \cdots \otimes e_{i_r} \otimes e^{j_1} \otimes \cdots \otimes e^{j_s}.
\]
\end{definition}

The second formula holds because both sides take the same value on every tuple $(e^{i_1}, \ldots, e^{i_r}, e_{j_1}, \ldots, e_{j_s})$ of basis vectors, and a multilinear map is
determined by these values. The components of a vector $v = \sum_iv^ie_i$ are its coordinates $v^i = e^i(v)$, with the index up; those of a linear form
$\varphi = \sum_j\varphi_je^j$ are $\varphi_j = \varphi(e_j)$, with the index down. The components of an endomorphism $A$, viewed as a tensor of type $(1,1)$, are
$A^i_j = e^i(Ae_j)$, the entries of its matrix (row $i$, column $j$); the components of a bilinear form $B$ are $B_{ij} = B(e_i, e_j)$, the entries of its Gram matrix
(Definition~\ref{linear-algebra:definition-bilinear-form}).

\begin{proposition}[Transformation law]\label{multilinear-algebra:proposition-tensor-transformation}
Let $(e^{\prime}_1, \ldots, e^{\prime}_n)$ be a second basis of $V$, with $e^{\prime}_j = \sum_iP^i_je_i$ for an invertible matrix $P = (P^i_j)$ (row $i$, column $j$), and let
$Q = (Q^i_j) = P^{-1}$.
\begin{enumerate}
\item The dual bases are related by $e^{\prime j} = \sum_iQ^j_ie^i$, and the coordinates of a vector by $v^{\prime i} = \sum_jQ^i_jv^j$.
\item The components of $T \in T^r_s(V)$ in the two bases are related by
\[
T^{\prime\, i_1 \ldots i_r}_{j_1 \ldots j_s} = \sum_{k_1, \ldots, k_r, l_1, \ldots, l_s} Q^{i_1}_{k_1} \cdots Q^{i_r}_{k_r}\, P^{l_1}_{j_1} \cdots P^{l_s}_{j_s}\; T^{k_1 \ldots k_r}_{l_1 \ldots l_s}.
\]
\item Conversely, suppose that to every basis $\mathbf{e}$ of $V$ there is assigned a family of numbers $T(\mathbf{e})^{i_1 \ldots i_r}_{j_1 \ldots j_s}$, such that the families
assigned to any two bases are related by the formula of (2). Then there is a unique $T \in T^r_s(V)$ whose components in every basis $\mathbf{e}$ are the numbers
$T(\mathbf{e})^{i_1 \ldots i_r}_{j_1 \ldots j_s}$.
\end{enumerate}
\end{proposition}

\begin{proof}
(1) The linear form $\sum_iQ^j_ie^i$ takes on $e^{\prime}_m$ the value $\sum_iQ^j_ie^i\bigl(\sum_lP^l_me_l\bigr) = \sum_iQ^j_iP^i_m = (QP)^j_m = \delta^j_m$, so it is the $j$-th vector of the
dual basis of $(e^{\prime}_m)$. Hence $v^{\prime i} = e^{\prime i}(v) = \sum_jQ^i_je^j(v) = \sum_jQ^i_jv^j$.

(2) By definition and (1),
\[
T^{\prime\, i_1 \ldots i_r}_{j_1 \ldots j_s} = T\Bigl(\sum_{k_1}Q^{i_1}_{k_1}e^{k_1}, \ldots, \sum_{k_r}Q^{i_r}_{k_r}e^{k_r}, \sum_{l_1}P^{l_1}_{j_1}e_{l_1}, \ldots, \sum_{l_s}P^{l_s}_{j_s}e_{l_s}\Bigr),
\]
and expanding by multilinearity gives the formula.

(3) Uniqueness holds because a tensor is determined by its components in one basis (Definition~\ref{multilinear-algebra:definition-tensor-components}). For existence, fix a
basis $\mathbf{e}$ and let $T$ be the tensor with components $T(\mathbf{e})^{i_1 \ldots i_r}_{j_1 \ldots j_s}$ in $\mathbf{e}$. For another basis $\mathbf{e}^{\prime}$, the components of $T$ in
$\mathbf{e}^{\prime}$ are obtained from those in $\mathbf{e}$ by the formula of (2), and so are the numbers $T(\mathbf{e}^{\prime})$ by hypothesis; so they agree.
\end{proof}

Part (3) is the justification of the physicists' definition, ``a tensor is something that transforms like a tensor''. Part (1) explains the names. The basis vectors
change by $P$; the coordinates of a vector change by the inverse matrix $Q$, against the basis, and are called \emph{contravariant}; the components of a linear form,
$\varphi^{\prime}_j = \varphi(e^{\prime}_j) = \sum_lP^l_j\varphi_l$, change by $P$, together with the basis, and are called \emph{covariant}. A tensor of type $(r,s)$ has $r$ indices
of the first kind, written above, and $s$ of the second kind, written below. (These traditional names clash with the language of categories: $V \mapsto V$ is a covariant functor and
$V \mapsto V^*$ a contravariant one, Example~\ref{categories:example-double-dual}, while vectors are called contravariant. The physicists' names refer to the components, not to the
spaces.)

\begin{example}\label{multilinear-algebra:example-tensor-transformation}
Let us compare the transformation laws of an endomorphism and of a bilinear form; in both cases the components form an $n \times n$ matrix.
\begin{enumerate}
\item For an endomorphism $A$, a tensor of type $(1,1)$, the formula of
Proposition~\ref{multilinear-algebra:proposition-tensor-transformation} reads $A^{\prime\,i}_j = \sum_{k,l}Q^i_kA^k_lP^l_j$, that is, $A^{\prime} = P^{-1}AP$: the matrix changes by
\emph{similarity}, as in Chapter~\ref{linear-algebra}. Similar matrices have the same eigenvalues, trace and determinant, so these are invariants of the endomorphism.
\item For a bilinear form $B$, a tensor of type $(0,2)$, the formula reads $B^{\prime}_{ij} = \sum_{k,l}P^k_iP^l_jB_{kl}$, that is, $B^{\prime} = P^TBP$: the Gram matrix changes by
\emph{congruence}, as in Definition~\ref{linear-algebra:definition-bilinear-form}. Congruent matrices need not have the same eigenvalues: for $B = I$ and $P = 2I$, $P^TBP = 4I$.
So ``the eigenvalues of a bilinear form'' make no sense; what survives for real symmetric forms is the signature (Theorem~\ref{linear-algebra:theorem-sylvester}).
\end{enumerate}
The same matrix can therefore mean two different things, and only the position of the indices tells which. With an inner product the two can be converted into each other
(Definition~\ref{multilinear-algebra:definition-raising-lowering} below), and if both bases are orthonormal for an inner product on a real vector space, then $P$ is orthogonal,
$P^{-1} = P^T$, and the two transformation laws coincide.
\end{example}

\begin{example}\label{multilinear-algebra:example-tensors}
\begin{enumerate}
\item The identity of $V$ is the tensor $\sum_ie_i \otimes e^i$ of type $(1,1)$; its components are $\delta^i_j$ in every basis, as the transformation law confirms:
$\sum_{k,l}Q^i_k\delta^k_lP^l_j = (QP)^i_j = \delta^i_j$. Hence $\sum_ie_i \otimes e^i$ does not depend on the basis, which is not obvious from the formula.
\item A bilinear map $\mu \colon V \times V \to V$, for instance the multiplication of a finite-dimensional algebra or the bracket of a Lie algebra, is a tensor of type $(1,2)$.
Its components $c^k_{ij} = e^k(\mu(e_i, e_j))$, defined by $\mu(e_i, e_j) = \sum_kc^k_{ij}e_k$, are the \emph{structure constants} of the algebra.
\item The determinant in the basis $(e_i)$, the alternating form $(w_1, \ldots, w_n) \mapsto \det(e^i(w_j))$, is a tensor of type $(0,n)$. More generally, alternating forms of degree $p$
are the tensors of type $(0,p)$ that change sign when two arguments are exchanged; Section~\ref{multilinear-algebra:section-alternating-forms} studies them.
\item In Riemannian geometry, the curvature of a Riemannian manifold assigns to each tangent space a tensor of type $(1,3)$, $(X, Y, Z) \mapsto R(X, Y)Z$, and its \emph{Ricci tensor}, of type $(0,2)$,
is obtained from it by a contraction (Chapter~\ref{riemannian-geometry}).
\end{enumerate}
\end{example}

The most important operation on tensors, after the tensor product, generalises the trace of an endomorphism: a vector slot and a covector slot are paired with each other.

\begin{definition}\label{multilinear-algebra:definition-contraction}
Let $r, s \geq 1$, $1 \leq a \leq r$ and $1 \leq b \leq s$. The \emph{contraction} of the $a$-th upper index with the $b$-th lower index is the linear map
$C^a_b \colon T^r_s(V) \to T^{r-1}_{s-1}(V)$ given by
\[
C^a_b(v_1 \otimes \cdots \otimes v_r \otimes \varphi^1 \otimes \cdots \otimes \varphi^s) = \varphi^b(v_a)\; v_1 \otimes \cdots \widehat{v_a} \cdots \otimes v_r \otimes \varphi^1 \otimes \cdots
\widehat{\varphi^b} \cdots \otimes \varphi^s,
\]
where the hat means that the factor is omitted.
\end{definition}

The map exists and is unique by Proposition~\ref{multilinear-algebra:proposition-tensor-basis}, since the right-hand side is multilinear in $(v_1, \ldots, v_r, \varphi^1, \ldots, \varphi^s)$.
Note that no basis is used in the definition.

\begin{proposition}\label{multilinear-algebra:proposition-contraction}
\begin{enumerate}
\item In components, contraction is summation over a repeated index:
\[
(C^a_bT)^{i_1 \ldots \widehat{i_a} \ldots i_r}_{j_1 \ldots \widehat{j_b} \ldots j_s} = \sum_{m=1}^nT^{i_1 \ldots m \ldots i_r}_{j_1 \ldots m \ldots j_s},
\]
with $m$ in the $a$-th upper and the $b$-th lower position.
\item For $r = s = 1$, $C^1_1 \colon \End(V) = T^1_1(V) \to k$ is the trace.
\item For $T \in T^1_2(V)$, viewed as a bilinear map $V \times V \to V$, the contraction $C^1_1T \in T^0_1(V) = V^*$ is the linear form $w \mapsto \operatorname{tr}\bigl(v \mapsto T(v, w)\bigr)$.
\end{enumerate}
\end{proposition}

\begin{proof}
(1) Both sides are linear in $T$, so it suffices to check the formula for a basis vector $T = e_{k_1} \otimes \cdots \otimes e_{k_r} \otimes e^{l_1} \otimes \cdots \otimes e^{l_s}$, whose components are
$T^{i_1 \ldots i_r}_{j_1 \ldots j_s} = \delta^{i_1}_{k_1} \cdots \delta^{i_r}_{k_r}\delta^{l_1}_{j_1} \cdots \delta^{l_s}_{j_s}$. By definition $C^a_bT = e^{l_b}(e_{k_a})$ times the tensor with the
$a$-th and $b$-th factors removed, and $e^{l_b}(e_{k_a}) = \delta^{l_b}_{k_a} = \sum_m\delta^m_{k_a}\delta^{l_b}_m$. Comparing components gives the formula.

(2) By (1), $C^1_1A = \sum_mA^m_m$, the sum of the diagonal entries of the matrix of $A$ (Definition~\ref{multilinear-algebra:definition-tensor-components}), which is the trace. Alternatively,
this is the last statement of Proposition~\ref{multilinear-algebra:proposition-hom-tensor}.

(3) By (1), $(C^1_1T)_j = \sum_mT^m_{mj}$. For fixed $w = e_j$ the endomorphism $v \mapsto T(v, e_j)$ has matrix entries $e^i(T(e_l, e_j)) = T^i_{lj}$, so its trace is $\sum_mT^m_{mj}$.
\end{proof}

\begin{remark}\label{multilinear-algebra:remark-einstein-convention}
Formulas with components contain many sums over indices that occur once above and once below, as in Propositions~\ref{multilinear-algebra:proposition-tensor-transformation}
and~\ref{multilinear-algebra:proposition-contraction}. Physicists and many geometers use the \emph{Einstein summation convention}: an index that appears once as an upper and once as a lower index is
summed over, and the summation sign is omitted. Thus $v = v^ie_i$, $A^{\prime\,i}_j = Q^i_kA^k_lP^l_j$, and the trace of $A$ is $A^i_i$. The convention has a built-in check: in a correct formula
every index that is not summed over appears in the same position on both sides, and each summed index appears exactly once up and once down. A formula such as $\sum_iv^iv^i$, which is not of this
form, is not invariant under change of basis; it becomes invariant only once a bilinear form is used to lower one of the indices, as in $\sum_{i,j}g_{ij}v^iv^j$. We do not use the convention in this book.
\end{remark}

Let $g$ be a nondegenerate symmetric bilinear form on $V$ (Definition~\ref{linear-algebra:definition-bilinear-form}), for instance an inner product when $k = \RR$. Its components
$g_{ij} = g(e_i, e_j)$ form an invertible symmetric matrix, and we write $g^{ij}$ for the entries of the inverse matrix, so that $\sum_jg^{ij}g_{jl} = \delta^i_l$.

\begin{definition}\label{multilinear-algebra:definition-raising-lowering}
The isomorphism $\flat \colon V \to V^*$, $v \mapsto v^\flat = g(v, -)$, and its inverse $\sharp \colon V^* \to V$, $\varphi \mapsto \varphi^\sharp$, are called \emph{lowering} and \emph{raising}
of indices. Applying $\flat$ to the $a$-th factor $V$ of $T^r_s(V)$, and moving the result to the last position, gives an isomorphism $T^r_s(V) \to T^{r-1}_{s+1}(V)$; applying $\sharp$ to a factor
$V^*$ gives an isomorphism $T^r_s(V) \to T^{r+1}_{s-1}(V)$. The \emph{trace with respect to $g$} of a tensor $B$ of type $(0,2)$ is
$\operatorname{tr}_gB = C^1_1(B^\sharp)$, the trace of the endomorphism obtained by raising the first index of $B$.
\end{definition}

\begin{proposition}\label{multilinear-algebra:proposition-raising-lowering}
\begin{enumerate}
\item In components, $(v^\flat)_i = \sum_jg_{ij}v^j$ and $(\varphi^\sharp)^i = \sum_jg^{ij}\varphi_j$.
\item For $B$ of type $(0,2)$, the endomorphism $B^\sharp$ is characterised by $g(B^\sharp w, v) = B(v, w)$ for all $v, w$, and $\operatorname{tr}_gB = \sum_{i,j}g^{ij}B_{ij}$.
\item If $k = \RR$, $g$ is an inner product and $(e_i)$ is orthonormal, then $g_{ij} = g^{ij} = \delta_{ij}$, raising and lowering do not change components, and
$\operatorname{tr}_gB = \sum_iB(e_i, e_i)$.
\end{enumerate}
\end{proposition}

\begin{proof}
(1) $(v^\flat)_i = v^\flat(e_i) = g(v, e_i) = \sum_jv^jg_{ji}$, and $g_{ji} = g_{ij}$. So $\flat$ has matrix $(g_{ij})$ in the bases $(e_i)$ and $(e^i)$, and its inverse $\sharp$ has matrix $(g^{ij})$.

(2) Write $B = \sum_{i,j}B_{ij}e^i \otimes e^j$. Raising the first index replaces $e^i$ by $(e^i)^\sharp = \sum_lg^{li}e_l$, so $B^\sharp = \sum_{i,j,l}g^{li}B_{ij}\,e_l \otimes e^j$, the endomorphism
$w \mapsto \sum_{i,j,l}g^{li}B_{ij}e^j(w)\,e_l = \sum_iB(e_i, w)(e^i)^\sharp$. Hence $g(B^\sharp w, v) = \sum_iB(e_i, w)e^i(v) = B(v, w)$; this characterises $B^\sharp w$ because $g$ is nondegenerate. By
Proposition~\ref{multilinear-algebra:proposition-contraction}(1) the trace is $\sum_m(B^\sharp)^m_m = \sum_{m,i}g^{mi}B_{im}$, which is $\sum_{i,j}g^{ij}B_{ij}$ since $g^{mi} = g^{im}$.

(3) The Gram matrix of an orthonormal basis is the identity, which is its own inverse; now apply (1) and (2).
\end{proof}

In particular $\operatorname{tr}_gB$ does not depend on the basis, being a contraction, although the formula $\sum_iB(e_i, e_i)$ is only valid for orthonormal bases. For $B = g$ one finds
$\operatorname{tr}_gg = \sum_{i,j}g^{ij}g_{ij} = n$. This is how the Ricci tensor and the scalar curvature of a Riemannian manifold are obtained from its curvature tensor
(Chapter~\ref{riemannian-geometry}), and how the Laplacian $\Delta f = \operatorname{tr}_g\operatorname{Hess}f$ is obtained from the Hessian.

\begin{exercise}\label{multilinear-algebra:exercise-tensor-symmetric}
Assume $\operatorname{char}k \neq 2$. Show that every tensor of type $(0,2)$ is uniquely the sum of a symmetric and an alternating one, and that the symmetric ones form a subspace of dimension
$n(n+1)/2$ and the alternating ones a subspace of dimension $n(n-1)/2$. Show that the symmetric ones are preserved by the transformation law, that is, the notion of symmetry does not depend on
the basis.
\end{exercise}

\begin{solution}
Write $B = \frac12(B + B^t) + \frac12(B - B^t)$, where $B^t(v, w) = B(w, v)$; the first summand is symmetric and the second alternating, because $(B - B^t)(v, v) = 0$. If $B$ is both symmetric and
alternating, then $0 = B(v + w, v + w) = B(v, w) + B(w, v) = 2B(v, w)$, so $B = 0$; this gives uniqueness. In components, symmetric means $B_{ij} = B_{ji}$, so a symmetric tensor is determined by
its components $B_{ij}$ with $i \leq j$, which are arbitrary: dimension $n + \binom n2 = n(n+1)/2$. Alternating means $B_{ij} = -B_{ji}$ and $B_{ii} = 0$, so the components with $i < j$
are free: dimension $\binom n2 = n(n-1)/2$; the two dimensions add up to $n^2$. Symmetry is defined by the condition $B(v, w) = B(w, v)$ for all $v, w$, which does not mention a basis; equivalently,
if $B^T = B$ then $(P^TBP)^T = P^TB^TP = P^TBP$.
\end{solution}

\begin{exercise}\label{multilinear-algebra:exercise-contraction-composition}
Let $A, B \in \End(V) = T^1_1(V)$, so that $A \otimes B \in T^2_2(V)$. Show that $C^2_1(A \otimes B) = A \circ B$ and that $C^1_1C^2_1(A \otimes B) = \operatorname{tr}(AB)$, while
$C^1_1C^1_1(A \otimes B) = \operatorname{tr}A\operatorname{tr}B$. Deduce that $\operatorname{tr}(AB) = \operatorname{tr}(BA)$.
\end{exercise}

\begin{solution}
By linearity we may take $A = v \otimes \varphi$ and $B = w \otimes \psi$, so $A \otimes B = v \otimes w \otimes \varphi \otimes \psi$, $A(x) = \varphi(x)v$ and $B(x) = \psi(x)w$. Then
$C^2_1(A \otimes B) = \varphi(w)\,v \otimes \psi$, the endomorphism $x \mapsto \varphi(w)\psi(x)v$; and $(A \circ B)(x) = A(\psi(x)w) = \psi(x)\varphi(w)v$, the same. Applying $C^1_1$, the trace
(Proposition~\ref{multilinear-algebra:proposition-contraction}(2)), gives $\operatorname{tr}(AB)$, which equals $\varphi(w)\psi(v)$. On the other hand $C^1_1(A \otimes B) = \varphi(v)\,w \otimes \psi$
and then $C^1_1$ gives $\varphi(v)\psi(w) = \operatorname{tr}A\operatorname{tr}B$. Finally $\operatorname{tr}(BA)$ is computed by the same formula with the roles exchanged, $\psi(v)\varphi(w)$,
which is the same number as $\operatorname{tr}(AB) = \varphi(w)\psi(v)$; both sides being bilinear in $(A, B)$, the identity holds for all $A$, $B$.
\end{solution}

\section{Alternating forms}\label{multilinear-algebra:section-alternating-forms}

In this section $V$ has finite dimension $n$.

\begin{proposition}\label{multilinear-algebra:proposition-dual-exterior}
The pairing $\Lambda^p(V^*) \times \Lambda^p(V) \to k$ determined by
$\langle \varphi_1 \wedge \cdots \wedge \varphi_p, v_1 \wedge \cdots \wedge v_p \rangle = \det(\varphi_i(v_j))$ is well defined and perfect. Hence
$\Lambda^p(V^*) \cong \Lambda^p(V)^*$, which by Proposition~\ref{multilinear-algebra:proposition-universal-algebras}
is the space $\mathrm{Alt}^p(V)$ of alternating $p$-linear forms on $V$.
\end{proposition}

\begin{proof}
The determinant is multilinear and alternating both in the $\varphi_i$ and in the $v_j$ (using
Proposition~\ref{linear-algebra:proposition-determinant-properties}(2)), so it induces a bilinear pairing.
For a basis $(e_i)$ with dual basis $(e_i^*)$, $\langle e_I^*, e_J \rangle = \delta_{IJ}$, so the pairing is perfect.
\end{proof}

\begin{notation}\label{multilinear-algebra:notation-forms}
We identify $\Lambda^p(V^*)$ with $\mathrm{Alt}^p(V)$ via
Proposition~\ref{multilinear-algebra:proposition-dual-exterior}: an element $\alpha \in \Lambda^p(V^*)$ is evaluated
on vectors as $\alpha(v_1, \ldots, v_p) = \langle \alpha, v_1 \wedge \cdots \wedge v_p \rangle$. With this convention,
$(\varphi_1 \wedge \cdots \wedge \varphi_p)(v_1, \ldots, v_p) = \det(\varphi_i(v_j))$, without a factor $1/p!$.
\end{notation}

\begin{example}\label{multilinear-algebra:example-forms-dictionary}
\emph{The three descriptions of an alternating form, explicitly.} Let $e_1, \ldots, e_n$ be a basis of $V$ with dual basis $e_1^*, \ldots, e_n^*$, and for $I = \{i_1 < \cdots < i_p\}$ write
$e_I = e_{i_1} \wedge \cdots \wedge e_{i_p}$ and $e_I^* = e_{i_1}^* \wedge \cdots \wedge e_{i_p}^*$.
\begin{enumerate}
\item \emph{From $\Lambda^p(V^*)$ to $\mathrm{Alt}^p(V)$.} An element $\alpha = \sum_Ia_Ie_I^*$ acts on vectors $v_j = \sum_iv_{ij}e_i$ by
$\alpha(v_1, \ldots, v_p) = \sum_Ia_I\det(v_{i_kj})_{k,j}$, the sum of the coefficients times the $p \times p$ minors, with rows $I$, of the matrix of coordinates of the $v_j$.
\item \emph{From $\mathrm{Alt}^p(V)$ to $\Lambda^p(V^*)$.} An alternating $p$-linear map $A$ is $\sum_IA(e_{i_1}, \ldots, e_{i_p})\,e_I^*$: the coefficients are the values on increasing tuples of basis vectors.
Summing over all tuples instead, $A = \frac1{p!}\sum_{i_1, \ldots, i_p}A(e_{i_1}, \ldots, e_{i_p})\,e_{i_1}^* \wedge \cdots \wedge e_{i_p}^*$, since each increasing tuple appears $p!$ times, with the same
sign in $A$ and in the wedge product.
\item \emph{As a linear form on $\Lambda^pV$.} $\alpha$ sends $e_J$ to $a_J$, by Proposition~\ref{multilinear-algebra:proposition-dual-exterior}; so $(e_I^*)$ is the dual basis of $(e_J)$.
\end{enumerate}
For instance, on $\RR^3$ let $\alpha = e_1^* \wedge e_2^* + 2e_2^* \wedge e_3^*$. As a bilinear map, $\alpha(v, w) = (v_1w_2 - v_2w_1) + 2(v_2w_3 - v_3w_2) = v^TAw$ with the skew-symmetric matrix $A$ having
$A_{12} = 1$, $A_{23} = 2$, $A_{13} = 0$; as a linear form on $\Lambda^2\RR^3$ it sends $e_1 \wedge e_2 \mapsto 1$, $e_2 \wedge e_3 \mapsto 2$, $e_1 \wedge e_3 \mapsto 0$; and under the identification
$\Lambda^2\RR^3 \cong \RR^3$ of the cross product ($e_2 \wedge e_3 \mapsto e_1$, $e_3 \wedge e_1 \mapsto e_2$, $e_1 \wedge e_2 \mapsto e_3$) it is the vector $(2, 0, 1)$, with
$\alpha(v, w) = \langle(2, 0, 1), v \times w\rangle$. For differential forms the same dictionary holds at each point with $e_i^* = dx^i$ and $e_i = \partial_i$: $\omega = \sum_I\omega_I\,dx^I$ with
$\omega_I = \omega(\partial_{i_1}, \ldots, \partial_{i_p})$.

\emph{Conventions.} Some books identify $\Lambda^p(V^*)$ with $\mathrm{Alt}^p(V)$ through the \emph{alternation} map $\mathrm{Alt}(T)(v_1, \ldots, v_p) = \frac1{p!}\sum_\sigma\operatorname{sgn}(\sigma)T(v_{\sigma(1)}, \ldots, v_{\sigma(p)})$,
which gives $(\varphi_1 \wedge \cdots \wedge \varphi_p)(v_1, \ldots, v_p) = \frac1{p!}\det(\varphi_i(v_j))$. In our convention $\varphi_1 \wedge \cdots \wedge \varphi_p = p!\,\mathrm{Alt}(\varphi_1 \otimes \cdots \otimes \varphi_p)$,
with no factor in the evaluation; for $p = 2$, $(\varphi \wedge \psi)(v, w) = \varphi(v)\psi(w) - \varphi(w)\psi(v)$ for us, and half of it in the other convention. Formulas such as the invariant formula for $d$
(Chapter~\ref{differential-forms}) change by such factors between books.
\end{example}

\begin{proposition}\label{multilinear-algebra:proposition-wedge-formula}
For $\alpha \in \Lambda^p(V^*)$ and $\beta \in \Lambda^q(V^*)$,
\[
(\alpha \wedge \beta)(v_1, \ldots, v_{p+q}) = \sum_{\sigma} \operatorname{sgn}(\sigma)\, \alpha(v_{\sigma(1)}, \ldots, v_{\sigma(p)})\,
\beta(v_{\sigma(p+1)}, \ldots, v_{\sigma(p+q)}),
\]
the sum running over the \emph{shuffles}: permutations $\sigma \in S_{p+q}$ with $\sigma(1) < \cdots < \sigma(p)$ and
$\sigma(p+1) < \cdots < \sigma(p+q)$.
\end{proposition}

\begin{proof}
Both sides are bilinear in $(\alpha, \beta)$, so we may take $\alpha = \varphi_1 \wedge \cdots \wedge \varphi_p$ and
$\beta = \varphi_{p+1} \wedge \cdots \wedge \varphi_{p+q}$. The left side is $\det M$ with $M = (\varphi_i(v_j))$. In the Leibniz
formula $\det M = \sum_\pi \operatorname{sgn}(\pi)\prod_j M_{\pi(j) j}$, group the permutations $\pi$ according to the set
$A = \{j : \pi(j) \leq p\}$ of columns sent to the first $p$ rows. For fixed $A = \{a_1 < \cdots < a_p\}$ with complement
$\{b_1 < \cdots < b_q\}$, let $\sigma$ be the shuffle with $\sigma(i) = a_i$ and $\sigma(p + i) = b_i$. Every such $\pi$ is
uniquely $\pi = (\pi_1 \times \pi_2) \circ \sigma^{-1}$ with $\pi_1 \in S_p$ acting on $\{1, \ldots, p\}$ and $\pi_2$ on
$\{p+1, \ldots, p+q\}$, and $\operatorname{sgn}\pi = \operatorname{sgn}\pi_1 \operatorname{sgn}\pi_2 \operatorname{sgn}\sigma$. Summing over $\pi_1$ and $\pi_2$
gives $\operatorname{sgn}(\sigma)\det(\varphi_i(v_{a_j}))\det(\varphi_{p+i}(v_{b_j}))$, which is the term of $\sigma$ on the right.
\end{proof}

\begin{definition}\label{multilinear-algebra:definition-interior-product}
For $v \in V$ and $\alpha \in \Lambda^p(V^*)$ with $p \geq 1$, the \emph{interior product} $\iota_v\alpha \in \Lambda^{p-1}(V^*)$ is
$(\iota_v\alpha)(v_2, \ldots, v_p) = \alpha(v, v_2, \ldots, v_p)$; for $p = 0$, $\iota_v\alpha = 0$.
\end{definition}

\begin{proposition}\label{multilinear-algebra:proposition-interior-antiderivation}
$\iota_v(\alpha \wedge \beta) = \iota_v\alpha \wedge \beta + (-1)^p\alpha \wedge \iota_v\beta$ for $\alpha \in \Lambda^p(V^*)$, and
$\iota_v \iota_w = -\iota_w\iota_v$.
\end{proposition}

\begin{proof}
Expanding the determinant along the first column (Laplace expansion) gives
$\iota_v(\varphi_1 \wedge \cdots \wedge \varphi_p) = \sum_{i=1}^p (-1)^{i-1}\varphi_i(v)\,\varphi_1 \wedge \cdots \widehat{\varphi_i} \cdots \wedge \varphi_p$.
Applying this formula to a product of decomposable forms and splitting the sum according to whether $i$ is
among the first $p$ factors gives the first identity. The second holds because $\alpha(w, v, \ldots) = -\alpha(v, w, \ldots)$.
\end{proof}

\begin{proposition}\label{multilinear-algebra:proposition-symplectic-determinant}
Suppose $\operatorname{char} k = 0$, and let $\omega \in \Lambda^2(V^*)$ be nondegenerate, $\dim V = 2m$. Then
$\omega^m = \omega \wedge \cdots \wedge \omega \neq 0$, and every $f \in \GL(V)$ with $\omega(fv, fw) = \omega(v, w)$ has $\det f = 1$.
\end{proposition}

\begin{proof}
With a symplectic basis (Theorem~\ref{linear-algebra:theorem-symplectic-basis}) and its dual basis,
$\omega = \sum_i e_i^* \wedge f_i^*$, as both sides take the same values on pairs of basis vectors. The $2$-forms
$e_i^* \wedge f_i^*$ commute and square to zero, so $\omega^m = m!\, e_1^* \wedge f_1^* \wedge \cdots \wedge e_m^* \wedge f_m^* \neq 0$.
The map $f^* \colon V^* \to V^*$, $\varphi \mapsto \varphi \circ f$, satisfies $(\Lambda^2 f^*)(\omega) = \omega$, hence
$(\Lambda^{2m}f^*)(\omega^m) = \omega^m$; but $\Lambda^{2m}f^*$ is multiplication by $\det f^* = \det f$
(Proposition~\ref{multilinear-algebra:proposition-exterior-properties}(5) and
Proposition~\ref{linear-algebra:proposition-duality}(3)).
\end{proof}

The same computation, applied to an arbitrary alternating form instead of a nondegenerate one, produces a square root of the determinant of a skew-symmetric matrix. In the rest
of this section $k$ is a field of characteristic $0$, $V = k^{2m}$ with standard basis $(e_i)$ and dual basis $(e_i^*)$. To a skew-symmetric matrix $A \in M_{2m}(k)$, that
is $A^T = -A$, attach the $2$-form
\[
\alpha_A = \sum_{i < j}A_{ij}\,e_i^* \wedge e_j^* = \frac12\sum_{i, j}A_{ij}\,e_i^* \wedge e_j^* \in \Lambda^2(V^*) .
\]
The two expressions agree because $A_{ii} = 0$ and $A_{ji}\,e_j^* \wedge e_i^* = A_{ij}\,e_i^* \wedge e_j^*$. As an alternating form (Notation~\ref{multilinear-algebra:notation-forms}),
$\alpha_A(v, w) = \sum_{i < j}A_{ij}(v_iw_j - v_jw_i) = \sum_{i, j}A_{ij}v_iw_j = v^TAw$. Since $\Lambda^{2m}(V^*)$ is one-dimensional, spanned by $e_1^* \wedge \cdots \wedge e_{2m}^*$
(Proposition~\ref{multilinear-algebra:proposition-exterior-properties}(3)), the following definition makes sense.

\begin{definition}\label{multilinear-algebra:definition-pfaffian}
The \emph{Pfaffian} of a skew-symmetric matrix $A \in M_{2m}(k)$ is the scalar $\mathrm{Pf}(A)$ with
\[
\alpha_A^m = \alpha_A \wedge \cdots \wedge \alpha_A = m!\,\mathrm{Pf}(A)\,e_1^* \wedge \cdots \wedge e_{2m}^* .
\]
\end{definition}

\begin{proposition}\label{multilinear-algebra:proposition-pfaffian}
Let $A \in M_{2m}(k)$ be skew-symmetric.
\begin{enumerate}
\item $\displaystyle\mathrm{Pf}(A) = \frac{1}{2^mm!}\sum_{\sigma \in S_{2m}}\operatorname{sgn}(\sigma)\,A_{\sigma(1)\sigma(2)}A_{\sigma(3)\sigma(4)} \cdots A_{\sigma(2m-1)\sigma(2m)}$. So $\mathrm{Pf}$ is a
homogeneous polynomial of degree $m$ in the entries, with $\mathrm{Pf}(\lambda A) = \lambda^m\mathrm{Pf}(A)$; for instance
\[
\mathrm{Pf}\begin{pmatrix} 0 & a\\ -a & 0\end{pmatrix} = a, \qquad \mathrm{Pf}(A) = A_{12}A_{34} - A_{13}A_{24} + A_{14}A_{23} \quad (m = 2).
\]
\item $\mathrm{Pf}(g^TAg) = \det(g)\,\mathrm{Pf}(A)$ for every $g \in M_{2m}(k)$. In particular $\mathrm{Pf}$ is invariant under conjugation by $\mathrm{SO}(2m)$, and changes sign under
conjugation by an orthogonal matrix of determinant $-1$.
\item If $A$ is block diagonal with blocks $\begin{pmatrix} 0 & a_i\\ -a_i & 0\end{pmatrix}$, $i = 1, \ldots, m$, then $\mathrm{Pf}(A) = a_1a_2 \cdots a_m$.
\item $\mathrm{Pf}(A)^2 = \det A$.
\end{enumerate}
\end{proposition}

\begin{proof}
(1) Expanding the $m$-th power of $\alpha_A = \frac12\sum_{i,j}A_{ij}e_i^* \wedge e_j^*$ gives
\[
\alpha_A^m = \frac{1}{2^m}\sum_{i_1, j_1, \ldots, i_m, j_m}A_{i_1j_1} \cdots A_{i_mj_m}\,e_{i_1}^* \wedge e_{j_1}^* \wedge \cdots \wedge e_{i_m}^* \wedge e_{j_m}^* .
\]
A wedge product of $1$-forms with a repeated factor vanishes, so only the terms in which $(i_1, j_1, \ldots, i_m, j_m) = (\sigma(1), \ldots, \sigma(2m))$ for a permutation $\sigma$
survive, and for them $e_{\sigma(1)}^* \wedge \cdots \wedge e_{\sigma(2m)}^* = \operatorname{sgn}(\sigma)\,e_1^* \wedge \cdots \wedge e_{2m}^*$ by
Proposition~\ref{multilinear-algebra:proposition-exterior-properties}(1). Comparing with the definition gives the formula. For $m = 1$ the two permutations give
$\frac12(A_{12} - A_{21}) = A_{12}$. In general the $2^mm!$ permutations that differ by reordering the $m$ pairs (an even permutation of positions) and by swapping
the two entries of some pairs (which changes both $\operatorname{sgn}\sigma$ and the sign of one factor) give the same term, so $\mathrm{Pf}(A)$ is the sum over the ways
of splitting $\{1, \ldots, 2m\}$ into $m$ pairs $\{a_l < b_l\}$ of $\pm\prod_lA_{a_lb_l}$, the sign being that of the permutation $(a_1, b_1, \ldots, a_m, b_m)$. For $m = 2$ the
three splittings $\{12, 34\}$, $\{13, 24\}$, $\{14, 23\}$ give the permutations $1234$, $1324$, $1423$, of signs $+1$, $-1$, $+1$.

(2) Let $g^* \colon V^* \to V^*$, $\varphi \mapsto \varphi \circ g$, where $g$ acts on $V = k^{2m}$ by matrix multiplication. Then $g^*e_k^* = \sum_ig_{ki}e_i^*$, because
$e_k^*(ge_i) = g_{ki}$. The induced map $\Lambda g^*$ on $\Lambda(V^*)$ is multiplicative, and
\[
(\Lambda^2g^*)(\alpha_A) = \frac12\sum_{k, l}A_{kl}\Bigl(\sum_ig_{ki}e_i^*\Bigr) \wedge \Bigl(\sum_jg_{lj}e_j^*\Bigr) = \frac12\sum_{i, j}(g^TAg)_{ij}\,e_i^* \wedge e_j^* = \alpha_{g^TAg},
\]
$g^TAg$ being again skew-symmetric. Taking $m$-th powers, $\alpha_{g^TAg}^m = (\Lambda^{2m}g^*)(\alpha_A^m) = m!\,\mathrm{Pf}(A)\det(g)\,e_1^* \wedge \cdots \wedge e_{2m}^*$, since
$\Lambda^{2m}g^*$ is multiplication by $\det g^* = \det g$ (Proposition~\ref{multilinear-algebra:proposition-exterior-properties}(5) and
Proposition~\ref{linear-algebra:proposition-duality}(3)). For $g$ orthogonal, $g^TAg = g^{-1}Ag$.

(3) Here $\alpha_A = \sum_ia_i\,e_{2i-1}^* \wedge e_{2i}^*$. The $2$-forms $e_{2i-1}^* \wedge e_{2i}^*$ commute and square to zero, so
$\alpha_A^m = m!\,a_1 \cdots a_m\,e_1^* \wedge \cdots \wedge e_{2m}^*$, as in the proof of Proposition~\ref{multilinear-algebra:proposition-symplectic-determinant}.

(4) Suppose first that $A$ is invertible. Then $\omega(v, w) = v^TAw$ is a nondegenerate alternating form, and by Theorem~\ref{linear-algebra:theorem-symplectic-basis} it has a
symplectic basis $e_1^{\prime}, \ldots, e_m^{\prime}, f_1^{\prime}, \ldots, f_m^{\prime}$. Let $g$ be the matrix with columns $e_1^{\prime}, f_1^{\prime}, e_2^{\prime}, f_2^{\prime}, \ldots, e_m^{\prime}, f_m^{\prime}$
(in this order). The entries of $g^TAg$ are the values of $\omega$ on pairs of these vectors, so $g^TAg = J$ is block diagonal with blocks
$\bigl(\begin{smallmatrix} 0 & 1\\ -1 & 0\end{smallmatrix}\bigr)$. By (3), $\mathrm{Pf}(J) = 1$, and $\det J = 1$. By (2), $\det(g)\mathrm{Pf}(A) = 1$, while $\det(g)^2\det(A) = \det J = 1$;
so $\mathrm{Pf}(A)^2 = \det(g)^{-2} = \det A$. If $A$ is not invertible, choose $v \neq 0$ with $Av = 0$ and an invertible $g$ with first column $v$. The first column of $g^TAg$ is
$g^TAv = 0$, and by skew-symmetry so is its first row; hence $\alpha_{g^TAg}$ is a $2$-form in $e_2^*, \ldots, e_{2m}^*$ only, and its $m$-th power, of degree $2m$ in $2m - 1$ linear
forms, vanishes. So $\det(g)\mathrm{Pf}(A) = \mathrm{Pf}(g^TAg) = 0$, and $\mathrm{Pf}(A) = 0 = \det A$.
\end{proof}

Since $\mathrm{Pf}$ is a polynomial with rational coefficients, the identity in (2) is a polynomial identity in the entries of $A$ and $g$; so it remains true when the entries
of $A$ are taken in any commutative $\QQ$-algebra, for instance the algebra of differential forms of even degree on a manifold. This is how the Pfaffian is applied to curvature
matrices in Chapter~\ref{characteristic-classes}, where it produces the Euler form.

\begin{example}\label{multilinear-algebra:example-pfaffian}
For the $4 \times 4$ skew-symmetric matrix with $A_{12} = 1$, $A_{13} = 2$, $A_{14} = 3$, $A_{23} = 4$, $A_{24} = 5$, $A_{34} = 6$, the formula gives
$\mathrm{Pf}(A) = 1 \cdot 6 - 2 \cdot 5 + 3 \cdot 4 = 8$, and indeed $\det A = 64$. A real skew-symmetric matrix has $\det A = \mathrm{Pf}(A)^2 \geq 0$; this explains why the
determinant of a skew-symmetric matrix of even size is a perfect square, while one of odd size is $0$ (as $\det A = \det A^T = \det(-A) = (-1)^n\det A$).
\end{example}

\section{Orientations and the Hodge star}\label{multilinear-algebra:section-hodge-star}

\begin{definition}\label{multilinear-algebra:definition-orientation}
Let $V$ be a real vector space of dimension $n \geq 1$. Two ordered bases are \emph{equally oriented} if the
change of basis matrix has positive determinant. This is an equivalence relation with two classes, the
\emph{orientations} of $V$; equivalently, an orientation is one of the two components
$\{c\,e_1 \wedge \cdots \wedge e_n : c > 0\}$, $\{c\,e_1 \wedge \cdots \wedge e_n : c < 0\}$ of $\Lambda^n V \setminus \{0\}$.
\end{definition}

\begin{proposition}\label{multilinear-algebra:proposition-induced-inner-product}
Let $(V, g)$ be a real inner product space of dimension $n$. There is a unique inner product on $\Lambda^p(V)$
with $\langle v_1 \wedge \cdots \wedge v_p, w_1 \wedge \cdots \wedge w_p \rangle = \det(g(v_i, w_j))$. If $(e_i)$ is orthonormal, the $e_I$
are orthonormal. If $V$ is oriented, the element $\mathrm{vol} = e_1 \wedge \cdots \wedge e_n$ does not depend on the choice
of an oriented orthonormal basis; it is the \emph{volume element}.
\end{proposition}

\begin{proof}
The formula is multilinear and alternating in the $v_i$ and in the $w_j$, so it induces a bilinear form, which is
symmetric by Proposition~\ref{linear-algebra:proposition-determinant-properties}(2). On an orthonormal basis
$\langle e_I, e_J \rangle = \delta_{IJ}$, so the form is positive definite. If $(f_i)$ is another oriented orthonormal basis,
$f_j = \sum_i a_{ij}e_i$ with $A \in \mathrm{O}(n)$ and $\det A > 0$, hence $\det A = 1$, and
$f_1 \wedge \cdots \wedge f_n = \det(A)\, e_1 \wedge \cdots \wedge e_n$.
\end{proof}

\begin{remark}\label{multilinear-algebra:remark-change-of-basis-top}
\emph{What a change of basis does to $e_1 \wedge \cdots \wedge e_n$.} Let $(e_1, \ldots, e_n)$ and $(f_1, \ldots, f_n)$ be two bases of $V$, and $A = (a_{ij})$ the change of basis matrix, $f_j = \sum_ia_{ij}e_i$, whose $j$-th
column contains the coordinates of $f_j$. Applying Proposition~\ref{multilinear-algebra:proposition-exterior-properties}(5) to the linear map $e_i \mapsto f_i$, whose matrix is $A$,
\[
f_1 \wedge \cdots \wedge f_n = \det(A)\,e_1 \wedge \cdots \wedge e_n .
\]
The dual bases change by the inverse transpose, $f^{*}_i = \sum_j(A^{-1})_{ij}\,e^{*}_j$, so $f^{*}_1 \wedge \cdots \wedge f^{*}_n = \det(A)^{-1}\,e^{*}_1 \wedge \cdots \wedge e^{*}_n$; and on $\Lambda^p$ the entries of $A$ are replaced by its
$p \times p$ minors (Proposition~\ref{multilinear-algebra:proposition-minors}). So the top exterior power sees a change of basis only through one number, $\det A$, and each condition one imposes on the bases is a
condition on the matrix:
\begin{enumerate}
\item \emph{Two arbitrary bases:} $A$ is invertible, $\det A \neq 0$, and $f_1 \wedge \cdots \wedge f_n$ is an arbitrary nonzero multiple of $e_1 \wedge \cdots \wedge e_n$. Every nonzero element of the line $\Lambda^nV$ arises from some basis.
\item \emph{Equally oriented bases:} $\det A > 0$ (Definition~\ref{multilinear-algebra:definition-orientation}), the matrices of the group $\mathrm{GL}_n^+(\RR)$. The two wedge products are positive multiples of each other, so they lie in the same
half-line of $\Lambda^nV \setminus \{0\}$: this is why an orientation can be described as such a half-line.
\item \emph{Two orthonormal bases} of an inner product space: $A^TA = I$, that is, $A \in \mathrm{O}(n)$, because the entries of $A^TA$ are the inner products $\langle f_j, f_k\rangle = \sum_ia_{ij}a_{ik}$. Then
$\det(A)^2 = \det(A^T)\det(A) = \det(I) = 1$, so $\det A = \pm1$ and $f_1 \wedge \cdots \wedge f_n = \pm e_1 \wedge \cdots \wedge e_n$.
\item \emph{Two oriented orthonormal bases:} both conditions, so $\det A = 1$, that is, $A \in \mathrm{SO}(n) = \{A \in \mathrm{O}(n) : \det A = 1\}$, and $f_1 \wedge \cdots \wedge f_n = e_1 \wedge \cdots \wedge e_n$. This is the
volume element of Proposition~\ref{multilinear-algebra:proposition-induced-inner-product}: it is the only top degree element that every oriented orthonormal basis produces.
\item \emph{The converse fails:} $\det A = 1$ does not force the bases to be orthonormal. For $f_1 = e_1$, $f_2 = e_1 + e_2$ in $\RR^2$, $A = \bigl(\begin{smallmatrix} 1 & 1\\ 0 & 1\end{smallmatrix}\bigr)$, a shear, and
$f_1 \wedge f_2 = e_1 \wedge e_1 + e_1 \wedge e_2 = e_1 \wedge e_2$ although $f_2$ is not a unit vector. The element $e_1 \wedge \cdots \wedge e_n$ remembers exactly an orientation and a volume: it determines the
basis up to the matrices of $\mathrm{SL}_n(\RR) = \{\det A = 1\}$, which preserve orientation and volume but not lengths and angles.
\end{enumerate}
The volume interpretation can be made precise with the inner product of Proposition~\ref{multilinear-algebra:proposition-induced-inner-product}. If $(e_i)$ is orthonormal, then
$|f_1 \wedge \cdots \wedge f_n|^2 = \det(\langle f_j, f_k\rangle)_{j,k} = \det(A^TA) = (\det A)^2$, so
\[
|f_1 \wedge \cdots \wedge f_n| = |\det A| = \sqrt{\det(\langle f_j, f_k\rangle)} ,
\]
the volume of the parallelepiped spanned by $f_1, \ldots, f_n$, computed from the Gram matrix of the vectors alone. The same holds for $p < n$ vectors: $|v_1 \wedge \cdots \wedge v_p| = \sqrt{\det(\langle v_j, v_k\rangle)}$ is the
$p$-dimensional volume of the parallelepiped they span, which is the formula behind the area of a parametrised surface. Over $\CC$, with a Hermitian inner product, two orthonormal bases differ by a unitary
matrix, and $|\det A| = 1$ but $\det A$ can be any complex number of modulus $1$; so a Hermitian inner product determines the element $e_1 \wedge \cdots \wedge e_n$ of the complex line $\Lambda^n_\CC V$ only up to a phase, and
choosing it is an extra datum, a fact behind
the holomorphic volume forms of Calabi--Yau manifolds (Chapter~\ref{calabi-yau}).
\end{remark}

\begin{theorem}\label{multilinear-algebra:theorem-hodge-star}
Let $(V, g)$ be an oriented real inner product space of dimension $n$. There is a unique linear map
$* \colon \Lambda^p(V) \to \Lambda^{n-p}(V)$ with
\[
\alpha \wedge *\beta = \langle \alpha, \beta \rangle\, \mathrm{vol} \qquad \text{for all } \alpha, \beta \in \Lambda^p(V).
\]
For an oriented orthonormal basis, $*e_I = \varepsilon_I e_{I^c}$ where $e_I \wedge e_{I^c} = \varepsilon_I\,\mathrm{vol}$,
$\varepsilon_I = \pm 1$. The map $*$ is an isometry, $** = (-1)^{p(n-p)}$ on $\Lambda^p(V)$, and
$\alpha \wedge *\beta = \beta \wedge *\alpha$.
\end{theorem}

\begin{proof}
The pairing $\Lambda^p(V) \times \Lambda^{n-p}(V) \to \Lambda^n(V) = \RR\,\mathrm{vol}$ is perfect, since $e_I \wedge e_J$ is
$\pm\mathrm{vol}$ if $J = I^c$ and $0$ otherwise. Hence for each $\beta$ there is a unique $*\beta$ representing the functional
$\alpha \mapsto \langle \alpha, \beta \rangle$, and it depends linearly on $\beta$. The formula: $e_J \wedge \varepsilon_I e_{I^c}$ is
$\varepsilon_I^2\,\mathrm{vol} = \mathrm{vol}$ if $J = I$ and $0$ otherwise, which is $\langle e_J, e_I \rangle\mathrm{vol}$. So $*$ maps the
orthonormal basis $(e_I)$ to $(\pm e_{I^c})$, hence is an isometry. Next,
$**e_I = \varepsilon_I\varepsilon_{I^c}e_I$, and $e_{I^c} \wedge e_I = (-1)^{p(n-p)}e_I \wedge e_{I^c}$ by
Proposition~\ref{multilinear-algebra:proposition-exterior-properties}(1), so
$\varepsilon_{I^c} = (-1)^{p(n-p)}\varepsilon_I$. The last identity follows from the symmetry of the inner product.
\end{proof}

\begin{example}\label{multilinear-algebra:example-star-r3}
On $\RR^3$ with the standard orientation, $*e_1 = e_2 \wedge e_3$, $*e_2 = e_3 \wedge e_1$, $*e_3 = e_1 \wedge e_2$,
$*1 = \mathrm{vol}$ and $*\mathrm{vol} = 1$. The map $(v, w) \mapsto *(v \wedge w)$ is the cross product. On $\RR^4$, $*$ is an
involution of $\Lambda^2$ (since $(-1)^{2 \cdot 2} = 1$), which splits into $\pm 1$-eigenspaces of dimension $3$: the self-dual
and anti-self-dual $2$-vectors.
\end{example}

The inner product of Proposition~\ref{multilinear-algebra:proposition-induced-inner-product} gives a norm $|\xi|$ on $p$-vectors and, through the pairing of Notation~\ref{multilinear-algebra:notation-forms}, a norm $|\varphi|$ on
$p$-forms. For the geometry of $p$-dimensional areas another pair of norms is better adapted: they treat the \emph{simple} $p$-vectors $v_1 \wedge \cdots \wedge v_p$, which represent oriented $p$-dimensional
parallelepipeds, as the basic objects.

\begin{definition}\label{multilinear-algebra:definition-comass}
Let $(V, g)$ be a real inner product space and $1 \leq p \leq \dim V$. A $p$-vector $\xi \in \Lambda^p(V)$ is \emph{simple} if $\xi = v_1 \wedge \cdots \wedge v_p$ for some $v_i \in V$. The \emph{comass} of a $p$-form
$\varphi \in \Lambda^p(V^*)$ is
\[
\|\varphi\| = \sup\{\varphi(\xi) : \xi \text{ simple}, |\xi| = 1\} = \sup\{\varphi(v_1, \ldots, v_p) : v_1, \ldots, v_p \text{ orthonormal}\},
\]
and the \emph{mass} of a $p$-vector $\xi$ is $\|\xi\| = \sup\{\varphi(\xi) : \|\varphi\| \leq 1\}$.
\end{definition}

The two descriptions of the comass agree: if $v_1, \ldots, v_p$ are linearly independent, Gram--Schmidt (Theorem~\ref{linear-algebra:theorem-gram-schmidt}) replaces them by an orthonormal family
$u_1, \ldots, u_p$ with $v_1 \wedge \cdots \wedge v_p = c\,u_1 \wedge \cdots \wedge u_p$, $c > 0$, since the change of basis is triangular with positive diagonal; and $|u_1 \wedge \cdots \wedge u_p| = 1$ by
Proposition~\ref{multilinear-algebra:proposition-induced-inner-product}. So the unit simple $p$-vectors are the $u_1 \wedge \cdots \wedge u_p$ with $(u_i)$ orthonormal: the oriented $p$-planes of $V$, each with its
unit area element.

\begin{proposition}\label{multilinear-algebra:proposition-comass}
Let $(V, g)$ be a real inner product space of dimension $n$, and $1 \leq p \leq n$.
\begin{enumerate}
\item The comass and the mass are norms, and $\|\varphi\| \leq |\varphi|$ and $|\xi| \leq \|\xi\|$ for all $\varphi$ and $\xi$.
\item For simple $\xi$, $\|\xi\| = |\xi|$. For $p = 1$, comass and mass equal the Euclidean norms.
\item If $\xi = \xi_1 + \cdots + \xi_m$ with simple $\xi_i$, then $\|\xi\| \leq |\xi_1| + \cdots + |\xi_m|$, and $\varphi(\xi) \leq \|\varphi\|\,\|\xi\|$ for all $\varphi$ and $\xi$.
\end{enumerate}
\end{proposition}

\begin{proof}
(1) By Cauchy--Schwarz, with $\varphi$ identified with an element of $\Lambda^p(V)$ via the inner product, $\varphi(\xi) \leq |\varphi|\,|\xi| = |\varphi|$ for unit $\xi$, so the supremum is finite and $\|\varphi\| \leq |\varphi|$.
The comass is clearly subadditive and homogeneous, and $\|\varphi\| \geq |\varphi(\xi)|$ for unit simple $\xi$, as $-\xi$ is also unit simple. So $\|\varphi\| = 0$ forces $\varphi = 0$, since the unit simple vectors include the orthonormal basis $(e_I)$ of
$\Lambda^p(V)$, on which $\varphi$ then vanishes. For $\xi = \sum_I\xi^Ie_I$ and $\|\varphi\| \leq 1$, $\varphi(\xi) = \sum_I\xi^I\varphi(e_I) \leq \sum_I|\xi^I|$, so $\|\xi\|$ is finite;
it is subadditive and homogeneous, being a supremum of linear functions of $\xi$; and $\|\xi\| \geq \varphi(\xi)$ for $\varphi = \xi/|\xi|$, viewed as a form, which has $\|\varphi\| \leq |\varphi| = 1$, so $\|\xi\| \geq |\xi|^2/|\xi| = |\xi|$.

(2) For simple $\xi \neq 0$, $\xi/|\xi|$ is a unit simple vector, so $\varphi(\xi) \leq \|\varphi\|\,|\xi|$, and taking the supremum over $\|\varphi\| \leq 1$ gives $\|\xi\| \leq |\xi|$; the reverse is (1). For $p = 1$ every
vector is simple, so both definitions reduce to the Euclidean ones.

(3) For $\|\varphi\| \leq 1$, $\varphi(\xi) = \sum_i\varphi(\xi_i) \leq \sum_i|\xi_i|$ by (2); take the supremum. The second inequality is the definition of the dual norm.
\end{proof}

In fact $\|\xi\|$ is the smallest value of $|\xi_1| + \cdots + |\xi_m|$ over all decompositions of $\xi$ into simple $p$-vectors \cite[1.8.1]{federer-gmt}: the mass of a $p$-vector is its total area, computed as
economically as possible from oriented parallelepipeds. By (2), the norms can differ from the Euclidean ones only because of non-simple $p$-vectors, and the next example shows that they do.

\begin{example}\label{multilinear-algebra:example-comass}
On $\RR^4$ with the standard inner product, let $\omega = e^1 \wedge e^2 + e^3 \wedge e^4$ and $\xi = e_1 \wedge e_2 + e_3 \wedge e_4$.
\begin{enumerate}
\item $|\omega| = |\xi| = \sqrt 2$, as $e_{12}$ and $e_{34}$ are orthonormal.
\item $\|\omega\| = 1$. Let $J(u^1, u^2, u^3, u^4) = (-u^2, u^1, -u^4, u^3)$, an isometry. Then $\omega(u, v) = u^1v^2 - u^2v^1 + u^3v^4 - u^4v^3 = \langle Ju, v\rangle \leq |Ju|\,|v| = 1$ for orthonormal $u, v$, with
equality for $v = Ju$, as for $(u, v) = (e_1, e_2)$. (This is the case $n = 2$ of the Wirtinger inequality of Chapter~\ref{kahler}: $\omega$ is the K\"ahler form of $\CC^2$, and the planes with $v = Ju$ are the complex
lines.)
\item $\|\xi\| = 2$: by (2), $\|\xi\| \geq \omega(\xi) = 2$, and by Proposition~\ref{multilinear-algebra:proposition-comass}(3), $\|\xi\| \leq |e_{12}| + |e_{34}| = 2$. In particular $\xi$ is not simple, by part (2) of the
proposition; directly, $\xi \wedge \xi = 2e_{1234} \neq 0$, while $v \wedge w \wedge v \wedge w = 0$.
\end{enumerate}
So the Euclidean norm underestimates the size of $\xi$, which consists of two unit squares in orthogonal planes and should have total area $2$, and overestimates the size of $\omega$, which measures every unit
square by at most $1$.
\end{example}

\section{Forms of type $(p,q)$}\label{multilinear-algebra:section-pq-forms}

Let $V$ be a real vector space of dimension $2n$ with a complex structure $J$ and $W = V^* = \Hom_\RR(V, \RR)$.
Then $W_\CC = W \otimes_\RR \CC$ is identified with $\Hom_\RR(V, \CC)$, and $\Lambda^k(W_\CC)$ (exterior power over $\CC$)
with $\Lambda^k(W) \otimes_\RR \CC$, the complex-valued alternating $k$-forms on $V$.

\begin{definition}\label{multilinear-algebra:definition-pq}
Let $W^{1,0} \subseteq W_\CC$ be the space of $\RR$-linear maps $\varphi \colon V \to \CC$ with $\varphi(Jv) = i\varphi(v)$, that is, the
$\CC$-linear maps $(V, J) \to \CC$, and $W^{0,1} = \overline{W^{1,0}}$, the maps with $\varphi(Jv) = -i\varphi(v)$. For
$p, q \geq 0$ let $\Lambda^{p,q} \subseteq \Lambda^{p+q}(W_\CC)$ be the span of the products $\alpha \wedge \beta$ with
$\alpha \in \Lambda^p W^{1,0}$ and $\beta \in \Lambda^q W^{0,1}$. Its elements are the forms of \emph{type} $(p, q)$.
\end{definition}

\begin{proposition}\label{multilinear-algebra:proposition-pq-decomposition}
\begin{enumerate}
\item $W_\CC = W^{1,0} \oplus W^{0,1}$, and both summands have complex dimension $n$.
\item $\Lambda^k(W_\CC) = \bigoplus_{p+q=k}\Lambda^{p,q}$, with $\dim_\CC \Lambda^{p,q} = \binom{n}{p}\binom{n}{q}$, and
$\overline{\Lambda^{p,q}} = \Lambda^{q,p}$.
\item A real $2$-form $\alpha \in \Lambda^2(W)$ lies in $\Lambda^{1,1}$ if and only if $\alpha(Jv, Jw) = \alpha(v, w)$ for all $v, w$.
\end{enumerate}
\end{proposition}

\begin{proof}
(1) The map $J^* \colon W \to W$, $\varphi \mapsto \varphi \circ J$, is a complex structure on $W$, and $W^{1,0}$, $W^{0,1}$ are the
$\pm i$-eigenspaces of its complexification; apply Proposition~\ref{linear-algebra:proposition-complex-structure}.
(2) Choose a basis $\zeta_1, \ldots, \zeta_n$ of $W^{1,0}$; then $\bar\zeta_1, \ldots, \bar\zeta_n$ is a basis of $W^{0,1}$ and together
they form a basis of $W_\CC$. By Proposition~\ref{multilinear-algebra:proposition-exterior-properties}(3) the products
$\zeta_I \wedge \bar\zeta_{I'}$ with $|I| + |I'| = k$ form a basis of $\Lambda^k(W_\CC)$ (after reordering factors, which only changes
signs), and those with $|I| = p$, $|I'| = q$ span $\Lambda^{p,q}$. Complex conjugation, extended to forms, maps
$\zeta_I \wedge \bar\zeta_{I'}$ to $\bar\zeta_I \wedge \zeta_{I'}$.
(3) The map $\alpha \mapsto J^*\alpha$, $(J^*\alpha)(v, w) = \alpha(Jv, Jw)$, equals $\Lambda^2$ of $J^*$ on $W_\CC$. On $\Lambda^{2,0}$ it acts by
$i \cdot i = -1$, on $\Lambda^{0,2}$ by $(-i)(-i) = -1$, and on $\Lambda^{1,1}$ by $i(-i) = 1$. So $\Lambda^{1,1}$ is exactly the fixed space of
$J^*$ on $\Lambda^2(W_\CC)$.
\end{proof}

\begin{example}\label{multilinear-algebra:example-standard-pq}
On $\RR^{2n} = \CC^n$ with the standard structure (Example~\ref{linear-algebra:example-standard-structure}) let
$dx_j, dy_j$ be the dual basis and $dz_j = dx_j + i\,dy_j$, $d\bar z_j = dx_j - i\,dy_j$. Then $dz_j \in W^{1,0}$, since
$dz_j(J\partial_{x_j}) = dz_j(\partial_{y_j}) = i = i\,dz_j(\partial_{x_j})$ and similarly on $\partial_{y_j}$. The fundamental form
of the standard inner product is
\[
\omega = \sum_j dx_j \wedge dy_j = \frac{i}{2}\sum_j dz_j \wedge d\bar z_j,
\]
because $dz_j \wedge d\bar z_j = -2i\,dx_j \wedge dy_j$. It is a real form of type $(1,1)$, in accordance with
Proposition~\ref{multilinear-algebra:proposition-pq-decomposition}(3) and
Proposition~\ref{linear-algebra:proposition-compatible-metric}.
\end{example}

\section{Exercises}\label{multilinear-algebra:section-exercises}

\begin{exercise}\label{multilinear-algebra:exercise-symmetric-dimension}
Show that $\dim S^d(k^n) = \binom{n + d - 1}{d}$.
\end{exercise}

\begin{exercise}\label{multilinear-algebra:exercise-divisibility}
Let $v \in V$ be nonzero and $\alpha \in \Lambda^p(V)$. Show that $v \wedge \alpha = 0$ if and only if $\alpha = v \wedge \beta$ for some
$\beta \in \Lambda^{p-1}(V)$.
\end{exercise}

\begin{exercise}\label{multilinear-algebra:exercise-plucker}
Let $\operatorname{char} k \neq 2$. Show that $\omega \in \Lambda^2(V)$ is decomposable, $\omega = v \wedge w$, if and only if
$\omega \wedge \omega = 0$.
\end{exercise}

\begin{exercise}\label{multilinear-algebra:exercise-star-norm}
In the situation of Theorem~\ref{multilinear-algebra:theorem-hodge-star}, show that
$\alpha \wedge *\alpha = \|\alpha\|^2\,\mathrm{vol}$, and compute $*$ on $\Lambda^1(\RR^2)$: it is the rotation by $90$ degrees.
\end{exercise}

\begin{solution}
The first identity is the definition with $\beta = \alpha$. On $\RR^2$, $*e_1 = e_2$ and $*e_2 = -e_1$, since
$e_1 \wedge e_2 = \mathrm{vol}$ and $e_2 \wedge (-e_1) = \mathrm{vol}$. This is the standard complex structure $J$.
\end{solution}

\begin{exercise}\label{multilinear-algebra:exercise-pq-real-dimension}
Let $V = \CC^n$. Show that the real $(k,k)$-forms, that is, $\alpha \in \Lambda^{k,k}$ with $\bar\alpha = \alpha$, form a real vector space
of dimension $\binom{n}{k}^2$.
\end{exercise}

\begin{exercise}\label{multilinear-algebra:exercise-weil-operator}
With notation as in Section~\ref{multilinear-algebra:section-pq-forms}, let $C$ act on $\Lambda^{p,q}$ as multiplication by
$i^{p-q}$. Show that $C$ commutes with complex conjugation, so preserves real forms, and that $C$ is the map
induced on $\Lambda(W_\CC)$ by $J^*$.
\end{exercise}

\begin{exercise}\label{multilinear-algebra:exercise-pfaffian-blocks}
Let $X \in M_m(k)$, $\operatorname{char}k = 0$, and $A = \begin{pmatrix} 0 & X\\ -X^T & 0\end{pmatrix} \in M_{2m}(k)$. Show that $\mathrm{Pf}(A) = (-1)^{m(m-1)/2}\det X$, and check the
case $m = 2$ against the formula of Proposition~\ref{multilinear-algebra:proposition-pfaffian}(1). Show also that $\mathrm{Pf}$ of a block diagonal matrix with skew-symmetric
blocks $A_1$, $A_2$ of even sizes is $\mathrm{Pf}(A_1)\mathrm{Pf}(A_2)$.
\end{exercise}

\begin{solution}
Here $\alpha_A = \sum_{i,j}X_{ij}\,e_i^* \wedge e_{m+j}^* = \sum_{i=1}^me_i^* \wedge \varphi_i$ with $\varphi_i = \sum_jX_{ij}e_{m+j}^*$. The $2$-forms $e_i^* \wedge \varphi_i$ commute and square
to zero, so $\alpha_A^m = m!\,(e_1^* \wedge \varphi_1) \wedge \cdots \wedge (e_m^* \wedge \varphi_m)$. Moving all the $e_i^*$ to the front, $\varphi_j$ must pass $e_{j+1}^*, \ldots, e_m^*$, which
gives the sign $(-1)^{(m-1) + (m-2) + \cdots + 0} = (-1)^{m(m-1)/2}$; and $\varphi_1 \wedge \cdots \wedge \varphi_m = \det(X)\,e_{m+1}^* \wedge \cdots \wedge e_{2m}^*$ by
Proposition~\ref{multilinear-algebra:proposition-exterior-properties}(5). So $\alpha_A^m = m!\,(-1)^{m(m-1)/2}\det(X)\,e_1^* \wedge \cdots \wedge e_{2m}^*$. For $m = 2$,
$A_{12} = A_{34} = 0$, $A_{13} = X_{11}$, $A_{24} = X_{22}$, $A_{14} = X_{12}$, $A_{23} = X_{21}$, and the formula gives $-X_{11}X_{22} + X_{12}X_{21} = -\det X$, in accordance with
$(-1)^{1} = -1$. For the block diagonal matrix, of sizes $2m_1$ and $2m_2$, $\alpha_A = \alpha_{A_1} + \alpha_{A_2}$ with $\alpha_{A_1}$, $\alpha_{A_2}$ commuting $2$-forms in disjoint sets of
variables. In $(\alpha_{A_1} + \alpha_{A_2})^{m_1 + m_2}$ only the term $\binom{m_1 + m_2}{m_1}\alpha_{A_1}^{m_1}\alpha_{A_2}^{m_2}$ survives, since a power of $\alpha_{A_i}$ above $m_i$ vanishes, and it equals
$\frac{(m_1 + m_2)!}{m_1!m_2!}\,m_1!\mathrm{Pf}(A_1)\,m_2!\mathrm{Pf}(A_2)\,e_1^* \wedge \cdots \wedge e_{2m}^*$.
\end{solution}
